Question 1
 * Total Viscous Drag Force (F):
   
 * Minimizing Viscous Drag:
   To minimize F, differentiate with respect to y and set to zero:
   
 * Relative Position (y):u
   
Question 2
 * Upstream Hydrostatic Force (F_1):
   * Density of upstream liquid: \rho_1 = 1.45 \times 1000 = 1450\text{ kg/m}^3
   * Gate height h_{gate} = 1.2\text{ m}, width w = 2.0\text{ m}, area A = 2.4\text{ m}^2.
   * Distance from free surface to centroid of gate: \bar{h}_1 = 1.5 + \frac{1.2}{2} = 2.1\text{ m}
   * Force: F_1 = \rho_1 g A \bar{h}_1 = (1450)(9.81)(2.4)(2.1) = 71,691.48\text{ N} \approx 71.69\text{ kN}
   * Center of pressure from top of upstream liquid (y_{p1}):
     
   * Height above the hinge at the bottom:
     
 * Downstream Hydrostatic Force (F_2):
   * Density of water: \rho_2 = 1000\text{ kg/m}^3
   * Centroid from water surface: \bar{h}_2 = \frac{1.2}{2} = 0.6\text{ m}
   * Force: F_2 = \rho_2 g A \bar{h}_2 = (1000)(9.81)(2.4)(0.6) = 14,126.4\text{ N} \approx 14.13\text{ kN}
   * Center of pressure from top of downstream water (y_{p2}):
     
   * Height above the hinge at the bottom:
     
 * Part (i): Resultant Force & Position:
   * Net Resultant Force:
     
   * Location of Resultant Force (Height above bottom hinge h_{pR}):
     
 * Part (ii): Force Required at the Top to Open Gate (F_{top}):
   Taking moments about the bottom hinge:
   
Question 3
 * Given:
   * Cylinder mass m = 10\text{ kg}, cross-sectional area A = 0.1\text{ m}^2, total weight W = 10 \times 9.81 = 98.1\text{ N}
   * Liquid A: S_A = 0.8 \implies \rho_A = 800\text{ kg/m}^3; submerged depth in A: h_A = 0.1\text{ m}
   * Liquid B: S_B = 1.0 \implies \rho_B = 1000\text{ kg/m}^3; submerged depth in B: h_B = 0.125\text{ m}
 * Gauge Pressure at the Cylinder Bottom:
   
 * Buoyant Force (F_B) & String Tension (T):
   * Total Upward Buoyant Force:
     
   * Vertical Equilibrium (\sum F_y = 0):
     
Question 4
Part (i): Curved Surface Forces
 * Parameters: Radius R = 1.2\text{ m}, Width W = 3.0\text{ m}, Depth to top of quadrant h_0 = 2.4\text{ m}.
 * Horizontal Component (F_H):
   * Projected vertical height H = 1.2\text{ m}; Area A_p = 1.2 \times 3.0 = 3.6\text{ m}^2.
   * Depth to centroid of projected area: \bar{h} = 2.4 + \frac{1.2}{2} = 3.0\text{ m}.
   * F_H = \rho g A_p \bar{h} = (1000)(9.81)(3.6)(3.0) = 105,948\text{ N} \approx 105.95\text{ kN}.
   * Line of Action (y_{pH} below surface):
     
 * Vertical Component (F_V):
   * Weight of water vertically above curved surface:
     
   * F_V = \rho g V = (1000)(9.81)(12.0329) \approx 118,042.8\text{ N} \approx 118.04\text{ kN}\text{ (acting downward)}.
   * Line of Action (x_{pV} from vertical boundary R-Q):
     
 * Resultant Force (F_R) & Angle (\theta):
   * Magnitude:
     
   * Direction with horizontal:
     
Part (ii): Pressure at Bottom of 6m Tank
 * Oil Layer: h_{oil} = 2\text{ m}, \rho_{oil} = 0.8 \times 1000 = 800\text{ kg/m}^3
   
 * Water Layer: h_{water} = 4\text{ m}, \rho_{water} = 1000\text{ kg/m}^3
   
